\documentclass[11pt, letterpaper]{article}
\input{../../homework.tex}
\usepackage{titlepageBU}
\title{Problem Set 8}
\courseID{PY 355}
\courseSection{A1}
\begin{document}
\makeproblem
\section*{Problem 1}
\begin{enumerate}
	\item 
		\[
			\begin{pmatrix}
			1 & a \\
			0 & 1
			\end{pmatrix}\begin{pmatrix}
			1 & b \\
			0 & 1
			\end{pmatrix}=\begin{pmatrix}
			1 & b+a \\
			0 & 1
			\end{pmatrix}
		.\]
	\item 
		\[
			\det \begin{pmatrix}
			1 & a \\
			0 & 1
			\end{pmatrix}=1
		.\]
		\[
			\begin{pmatrix}
			1 & a \\
			0 & 1
			\end{pmatrix}^{-1} =\frac{1}{1} \begin{pmatrix}
			1 & -a \\
			-0 & 1
			\end{pmatrix}=\begin{pmatrix}
			1 & -a \\
			0 & 1
			\end{pmatrix}
		.\]
	\item 
		\[
			\left[ M_a,M_b \right] =M_aM_b-M_bM_a=\begin{pmatrix}
			1 & a+b \\
			0 & 1
			\end{pmatrix}-\begin{pmatrix}
			1 & b+a \\
			0 & 1
			\end{pmatrix}=0
		.\]
	\item Inducting on part 1 gives
		\[
			\prod_{n=0}^{N} M_{2^n} =\begin{pmatrix}
			1 & \sum_{n=0}^{N} 2^n \\
			0 & 1
			\end{pmatrix}
		.\]
		The series is geometric, and finite geometric series satisfy
		\[
			\sum_{n=1}^{N} ar^n=a\left( \frac{1-r^N}{1-r}  \right) 
		.\]
		It follows that
		\[
			\sum_{n=0}^{N} 2^n =\sum_{n=1}^{N+1} 2^{n-1}  =\frac{1}{2} \sum_{n=1}^{N+1} 2^{n}  =\frac{1}{2} \cdot 2\left( \frac{1-2^{N+1} }{1-2}  \right) =2^{N+1} -1
		.\]
		Therefore,
		\[
			\prod_{n=0}^{N} M_{2^n} =\begin{pmatrix}
			1 & 2^{N+1} -1 \\
			0 & 1
			\end{pmatrix}
		.\]
\end{enumerate}

\section*{Problem 2}
\begin{enumerate}
	\item I have absolutely no idea what the $\epsilon $ is supposed to be.

		ok a quick google search shows that $\epsilon_{ijk} $ is the Levi-Civita symbol, which can be intuitively understood to behave like a wedge product or an alternating tensor, in that $\epsilon_{ijk} =1$ if $ijk=123$, and is alternating in that each transposition induces a factor of $-1$, and this anticommutativity necessarily implies $\epsilon_{ijk} =0$ if $i=j$, $j=k$, or $k=i$. Due to this, the identity claims that $\sigma_i\sigma_i=0+\delta_{ii}1=1 $, so we verify this:
		\[
			\begin{pmatrix}
			0 & 1 \\
			1 & 0
			\end{pmatrix}\begin{pmatrix}
			0 & 1 \\
			1 & 0
			\end{pmatrix}=\begin{pmatrix}
			1 & 0 \\
			0 & 1
			\end{pmatrix},
		\]
		\[
			\begin{pmatrix}
			0 & -i \\
			i & 0
			\end{pmatrix}\begin{pmatrix}
			0 & -i \\
			i & 0
			\end{pmatrix}=\begin{pmatrix}
			-i^2 & 0 \\
			0 & -i^2
			\end{pmatrix}=\begin{pmatrix}
			1 & 0 \\
			0 & 1
			\end{pmatrix}
		,\]
		\[
			\begin{pmatrix}
			1 & 0 \\
			0 & -1
			\end{pmatrix}\begin{pmatrix}
			1 & 0 \\
			0 & -1
			\end{pmatrix}=\begin{pmatrix}
			1 & 0 \\
			0 & 1
			\end{pmatrix}
		.\]
		This verifies the identity for $i=j$. Moreover, we now have
		\[
			\sigma_i\sigma_j+\sigma_j\sigma _i=2\delta_{ij} 1
		.\]
		Note that
		\begin{align*}
			\sigma _i\sigma _j&=\sigma _i\sigma _j+\frac{1}{2} \sigma _j\sigma _i-\frac{1}{2} \sigma _j\sigma _i\\
					  &=\frac{1}{2} \left( \sigma _i\sigma _j-\sigma _j\sigma _i \right) +\frac{1}{2} \left( \sigma _i\sigma _j+\sigma _j\sigma _i \right)\\
					  &=\frac{1}{2} \left[ \sigma _i,\sigma _j \right] +\frac{1}{2} 2\delta_{ij} 1\\
					  &=\frac{1}{2} \left[ \sigma _i,\sigma _j \right] +\delta_{ij} 1
		.\end{align*}
		It thus suffices to show $1/2 \left[ \sigma _i,\sigma _j \right] =ie_{ijk} \sigma _k$. If we can show that $\left[ \sigma_1,\sigma _2 \right] $, $\left[ \sigma _2,\sigma _3 \right] $, and $\left[ \sigma _3,\sigma _1 \right] $ are equal to $2i\sigma_k$ for $k$ the indicie not appearing, then this suffices as a proof, since anticommutivity of the commutator is trivially equivalent to the alternating nature of the Levi-Citiva symbol. These computations are simple and can be brute forced.
		\[
			\left[ \sigma _1,\sigma _2 \right] =\begin{pmatrix}
			0 & 1 \\
			1 & 0
			\end{pmatrix}\begin{pmatrix}
			0 & -i \\
			i & 0
			\end{pmatrix}-\begin{pmatrix}
			0 & -i \\
			i & 0
			\end{pmatrix}\begin{pmatrix}
			0 & 1 \\
			1 & 0
			\end{pmatrix}=\begin{pmatrix}
			i & 0 \\
			0 & -i
			\end{pmatrix}-\begin{pmatrix}
			-i & 0 \\
			0 & i
			\end{pmatrix}=2i\sigma_3
		.\]
		\[
			\left[ \sigma _2,\sigma _3 \right] =\begin{pmatrix}
			0 & -i \\
			i & 0
			\end{pmatrix}\begin{pmatrix}
			1 & 0 \\
			0 & -1
			\end{pmatrix}-\begin{pmatrix}
			1 & 0 \\
			0 & -1
			\end{pmatrix}\begin{pmatrix}
			0 & -i \\
			i & 0
			\end{pmatrix}=\begin{pmatrix}
			0 & i \\
			i & 0
			\end{pmatrix}-\begin{pmatrix}
			0 & -i \\
			-i & 0
			\end{pmatrix}=2i\sigma_1
		.\]
		\[
			\left[ \sigma _3,\sigma _1 \right] =\begin{pmatrix}
			1 & 0 \\
			0 & -1
			\end{pmatrix}\begin{pmatrix}
			0 & 1 \\
			1 & 0
			\end{pmatrix}-\begin{pmatrix}
			0 & 1 \\
			1 & 0
			\end{pmatrix}\begin{pmatrix}
			1 & 0 \\
			0 & -1
			\end{pmatrix}=\begin{pmatrix}
			0 & 1 \\
			-1 & 0
			\end{pmatrix}-\begin{pmatrix}
			0 & -1 \\
			1 & 0
			\end{pmatrix}=\begin{pmatrix}
			0 & 2 \\
			-2 & 0
			\end{pmatrix}=2i\sigma_2
		.\]
		Therefore,
		\[
			\sigma _i\sigma _j=i\epsilon_{ijk} \sigma _k+\delta_{ij} 1
		.\]
	\item We have
		\[
			\left( \sigma\cdot p \right) ^2=\left( \sigma_i p_i \right) ^2=\sigma _ip_i\sigma _jp_j=p_ip_j\sigma _i\sigma _j
		.\]
		Using our previous result, we have
		\[
			p_ip_j\left( i \epsilon_{ijk}\sigma_k+\delta_{ij} 1 \right) 
		.\]
		Evaluating the term on the right, we have
		\[
			p_ip_j\delta_{ij}1 =p_ip_i1=p_i^21=\left| p \right| ^21
		.\]
		Evaluating the term on the left, we have
		\[
			ip_ip_j\epsilon_{ijk} \sigma _k
		.\]
		Note that for each term $(i,j)$, there is a corresponding term $(j,i)$, and the anticommutivity induced by the alternating Levi-Citiva term causes these off-diagonal terms to vanish. Similarly, by anticommutivity once again, the diagonal terms also vanish, and thus
		\[
			ip_ip_j\epsilon_{ijk} \sigma _k=0
		.\]
		Therefore,
		\[
			\left( \sigma \cdot p \right) =\left| p \right| ^21
		.\]
\end{enumerate}

\section*{Problem 3}

Let $A^2=0$. Then
\[
	\begin{pmatrix}
	a & b \\
	c & d
	\end{pmatrix}\begin{pmatrix}
	a & b \\
	c & d
	\end{pmatrix}=\begin{pmatrix}
	a^2+bc & ab+bd \\
	ac+cd & bc+d^2
	\end{pmatrix}=0
.\]
If $b=c=0$, then clearly $A=0$, so we can without loss of generality ignore this case. We have
\[
	b(a+d)=c(a+d)=0
.\]
Since we take at least one of $\left\{ b,c \right\} $ to be nonzero, we must have $a+d=0$. The matrix becomes
\[
	A=\begin{pmatrix}
	a & b \\
	c & -a
	\end{pmatrix},\qquad A^2=\begin{pmatrix}
	a^2+bc & ab-ab \\
	ac-ac & a^2+bc
	\end{pmatrix}=\begin{pmatrix}
	a^2+bc & 0 \\
	0 & a^2+bc
	\end{pmatrix}
.\]
Now we re-parameterize with $x$ and $y$. Choose $y\in\mathbb{C} $ such that $b=y^2$. Define $x\coloneqq a/y$, to enforce that $a=xy$. We then have
\[
	\left( xy \right)^2+y^2c=y^2\left( x^2+c \right) =0
.\]
It follows that $c=-x^2$, since $b\neq 0$ by assumption. Therefore, our matrix becomes
\[
	A=\begin{pmatrix}
	xy & y^2 \\
	-x^2 & -xy
	\end{pmatrix}
.\]
This differs from the matrix in the problem only in notation.

\section*{Problem 4}
\begin{enumerate}
	\item 
		\[
			\det \begin{pmatrix}
			1 & 1 & 1 \\
			0 & 1 & 0 \\
			1 & 0 & 0
			\end{pmatrix}=\det \begin{pmatrix}
			0 & 1 \\
			1 & 0
			\end{pmatrix}=-1
		.\]
	\item 
		\[
			\det \begin{pmatrix}
			1 & 2 \\
			3 & 1
			\end{pmatrix}=1^2-6=-5
		.\]
	\item 
		\[
			\det \begin{pmatrix}
			0 & \sqrt{2}  & 0 & 0 \\
			\sqrt{2}  & 0 & 3 & 0 \\
			0 & 3 & 0 & \sqrt{2}  \\
			0 & 0 & \sqrt{2}  & 0
			\end{pmatrix}=-\det \begin{pmatrix}
			\sqrt{2}  & 3 & 0 \\
			0 & 0 & \sqrt{2}  \\
			0 & \sqrt{2}  & 0
			\end{pmatrix}=\sqrt{2} \det \begin{pmatrix}
			0 & \sqrt{2}  \\
			\sqrt{2}  & 0
			\end{pmatrix}=2\sqrt{2} 
		.\]
\end{enumerate}

\section*{Problem 5}
Determinants of triangular matrices are simply the product of the diagonal entries, and $\det AB=\det A \det B$. We have
\[
	\det A=\left( 14\cdot 3 \right) \left( -8 \right) =-42\cdot 8=-336
.\]
\end{document}
